\scanpage{174}
\mspage{123}
% archive sheet number: 174
pour des catégories très particulières -- les seules que je connaisse sont les catégories discrètes (qu'il faut peut-être $=$ supposer finies en plus).

Notons qu'on a un diagramme de foncteurs canoniques

\diagram{\(\unsure{(*)|(\alpha)|(**)}\qquad\begin{tikzcd}[ampersand replacement=\&,column sep=small,row sep=small]
\& X_0\arrow[dl,"\alpha_1"']\arrow[dr,"\alpha_2"] \&\\
X_1\arrow[dr,"{\substack{i_1\\ \deleted{\unsure{\alpha_2}}}}"'] \& \substack{\text{non}\\\text{commuta-}\\\text{tif !}} \& X_2\arrow[dl,"{\deleted{\unsure{\alpha_1}}\,i_2}"]\\
\& \unsure{\mathscr{X}|\mathcal X} \& {,}
\end{tikzcd}\)}
lequel diagramme n'a rien de commutatif dans $\mathrm{Cat}$. \uline{Il est cependant commutatif dans} $\mathrm{Hot}_W$ (quel que soit le localiseur fondamental $W$), \inserted{et même dans $\unsure{\mathrm{Cat}_{\mathrm{hot}}|\mathrm{Cat}\,\mathrm{hot}}$,} car si on désigne par

\[
i_0:X_0\longrightarrow\mathscr{X}
\]
l'inclusion canonique, on a des \unsure{hom|homs} canoniques

\diagram{\begin{tikzcd}[ampersand replacement=\&]
i_1\alpha_{\unsure{1|a}}\arrow[r] \& i_0 \& i_2\alpha_2\arrow[l]
\end{tikzcd}}
ce qui montre que $i_1\alpha_1$ et $i_2\alpha_2$ sont homotopes, donc définissent la $=$ flèche dans $\mathrm{Hot}_W$.
